\documentclass[12pt]{article}
\usepackage[utf8]{inputenc}
\usepackage[T1]{fontenc}
\usepackage{hyperref}
\title{Wurrugu Language Resources for K-12 Curriculum Development}
\author{Compiled by Claude Sonnet 5.5 (claude-sonnet-5-5) for Andr\'as Kornai}
\date{}
\begin{document}
\maketitle

\section{Introduction}
Wurrugu (Wurango, Popham Bay language; ISO 639-3 wur) is an extinct Aboriginal language of the Cobourg Peninsula in the Northern Territory, Australia. It is known only from a few 19th-century wordlists and one rememberer. It is grouped with Marrgu in the proposed Marrku--Wurrugu family, and its classification is uncertain.

\section{Primary Educational Resources}
None exist. The language has no living speakers, and no teaching materials turned up in the searches.

\section{Linguistic and Reference Materials}
Overviews are at \url{https://en.wikipedia.org/wiki/Wurrugu_language} and \url{https://en.wikipedia.org/wiki/Marrku%E2%80%93Wurrugu_languages}. The AIATSIS Austlang record is at \url{https://collection.aiatsis.gov.au/austlang/language/N37}. The Glottolog family page is at \url{https://glottolog.org/resource/languoid/id/marr1257}.

\section{Religious and Community Resources}
None found.

\section{Digital and Audio Resources}
None found.

\section{Curriculum Development Context}
No K-12 curriculum in Wurrugu is possible with the available record. Any educational use would be limited to cultural and historical content about the Cobourg Peninsula, developed with the local Aboriginal communities and consulting the AIATSIS record.
\end{document}